\documentclass[conference]{IEEEtran}
\IEEEoverridecommandlockouts
% The preceding line is only needed to identify funding in the first footnote. If that is unneeded, please comment it out.
\usepackage{cite}
\usepackage{amsmath,amssymb,amsfonts}
\usepackage{algorithmic}
\usepackage{graphicx}
\usepackage{textcomp}
\usepackage{xcolor}
\def\BibTeX{{\rm B\kern-.05em{\sc i\kern-.025em b}\kern-.08em
    T\kern-.1667em\lower.7ex\hbox{E}\kern-.125emX}}
\usepackage{algorithm}
\usepackage{array}
%\usepackage[caption=false,font=normalsize,labelfont=sf,textfont=sf]{subfig}
\usepackage{stfloats}
\usepackage{url}
\usepackage{verbatim}
\usepackage{booktabs}
\hyphenation{op-tical net-works semi-conduc-tor IEEE-Xplore}
%\usepackage[bottom=1.5in]{geometry}
% updated with editorial comments 8/9/2021

\begin{document}

\title{Latent Diffusion Synthetic Data for Low-Data IMS Classification\thanks{DISTRIBUTION STATEMENT A. Approved for Public Release: distribution is unlimited. Cleared for Public Release.}}


\author{
\IEEEauthorblockN{1\textsuperscript{st} Kevin Metzler}
\IEEEauthorblockA{\textit{Mathematical Sciences}\\
\textit{Worcester Polytechnic Institute}\\
Worcester, MA, USA \\
kjmetzler@wpi.edu}
\and
\IEEEauthorblockN{2\textsuperscript{nd} Hy Lam}
\IEEEauthorblockA{\textit{Mathematical Sciences}\\
\textit{Worcester Polytechnic Institute}\\Worcester, MA, USA \\
hlam@wpi.edu}
\and
\IEEEauthorblockN{3\textsuperscript{rd} Cate Dunham}
\IEEEauthorblockA{\textit{Data Science}\\
\textit{Worcester Polytechnic Institute}\\Worcester, MA, USA \\
cmdunham@wpi.edu}
\and
%\IEEEauthorblockN{4\textsuperscript{th} Steven Tate}
%\IEEEauthorblockA{\textit{Worcester Polytechnic Institute}\\
%Worcester, MA, USA \\
%srtate@wpi.edu}
%\and
%\IEEEauthorblockN{4\textsuperscript{th} Zahra Bolkheiri}
%\IEEEauthorblockA{\textit{Paffenroth Research Group} \\
%\textit{Worcester Polytechnic Institute}\\
%Worcester, MA, USA \\
%zahrabolkhairi@gmail.com}
%\and
\IEEEauthorblockN{4\textsuperscript{th} Patrick C. Riley}
\IEEEauthorblockA{\textit{U.S. Army DEVCOM Chemical Biological Center} \\
Aberdeen Proving Ground, MD, USA \\
patrick.c.riley15.civ@army.mil}
\and
\IEEEauthorblockN{5\textsuperscript{th} Samir V. Deshpande}
\IEEEauthorblockA{\textit{U.S. Army DEVCOM Chemical Biological Center}\\
Aberdeen Proving Ground, MD, USA\\
samir.v.deshpande.civ@army.mil}
%\and
%\IEEEauthorblockN{6\textsuperscript{th} Casandra}
%\IEEEauthorblockA{\textit{CBC}\\
%, USA \\
%}
\and
\IEEEauthorblockN{6\textsuperscript{th} Thomas Cao}
\IEEEauthorblockA{\textit{Research Scientist}\\
\textit{Systems Planning \& Analysis}\\
\textit{DTRA RD-CBEI A\&AS}\\
Fort Belvoir, VA, USA \\
tnstcc@gmail.com}
\and
%\IEEEauthorblockN{4\textsuperscript{th} \textbf{FIXME} Edward Argenta}
%\IEEEauthorblockA{\textit{Interdisciplinary Physical Scientist}\\
%\textit{DTRA RD-CBEI S\&T Manager}\\
%Fort Belvoir, VA, USA \\
%chiawei.tsai.civ@mail.mil}
%\and
\IEEEauthorblockN{7\textsuperscript{th} Joshua Uzarski}
\IEEEauthorblockA{\textit{Research Chemist}\\
\textit{US Army}\\
\textit{DEVCOM Soldier Center}\\
\textit{Soldier Protection Division}\\
Natick, MA, USA \\
joshua.r.uzarski.civ@army.mil}
\and
\IEEEauthorblockN{8\textsuperscript{th} Randy Paffenroth}
\IEEEauthorblockA{\textit{Mathematical Sciences,}\\
\textit{Data Science, and Computer Science}\\
\textit{Worcester Polytechnic Institute}\\Worcester, MA, USA \\
rcpaffenroth@wpi.edu}
}
\IEEEpubidadjcol


\maketitle

\noindent\textit{Technical Area:} D. Adaptive Systems, Machine Learning, and Data Analytics \\
\textit{Topics:} 9. Self- and Semi-supervised Learning, 10. Deep Learning

\begin{abstract}
%In ML classifiers for ion mobility spectrometry (IMS), classifier performance is often limited by scarce labeled spectra. This extended summary reports preliminary findings on whether synthetic samples generated by our geometry-preserving latent diffusion model improve downstream classification. Using IMS spectra from 8 chemical compounds, we evaluate Random Forest(RF) and Multi-Layer Perceptron(MLP) classifiers trained with varying real/synthetic mixtures across 50 low-data regimes at 3 levels of synthetic sampling variance. Initial results show synthetic data provides consistent but modest gains in absolute accuracy(aa) for RFs (+0.0061aa to +0.0097aa), and large gains in aa for MLPs when synthetic variance is increased (std=1.5: +0.1519aa, std=2.0: +0.1304aa). At std=1.5, the untuned MLPs improve over real-only training in 49/50 regimes, with synthetic-only training outperforming real-only in 46/50 regimes. These preliminary results show that our manifold-aware synthetic generation can materially improve low-data classifier training, especially for high-capacity neural classifiers. Ongoing work that would be included in a full paper includes hyperparameter optimization for both classifiers and the generative model, evaluation of additional classifier architectures, and refinement of the diffusion model's loss function to further improve synthetic data quality.

Ion mobility spectrometry (IMS) classifier performance is often limited by scarce labeled spectra. This extended summary reports preliminary findings on whether synthetic samples generated by our geometry-preserving latent diffusion model improve downstream classification. Using IMS spectra from 8 chemical compounds, we evaluate Random Forest (RF) and Multi-Layer Perceptron (MLP) classifiers trained with varying real/synthetic mixtures across 50 low-data regimes at 3 levels of synthetic sampling variance. Results show synthetic data provides consistent but modest gains in absolute accuracy (aa) for RFs (+0.0061aa to +0.0097aa), and large gains for MLPs when synthetic variance is increased (std=1.5: +0.1519aa, std=2.0: +0.1304aa). At std=1.5, untuned MLPs improve over real-only training in 49/50 regimes, with synthetic-only training outperforming real-only in 46/50 regimes. These preliminary results show that our manifold-aware synthetic generation can materially improve low-data classifier training, especially for high-capacity neural classifiers. Ongoing work includes hyperparameter optimization for the classifiers and the generative model, evaluation of additional classifier architectures, and refinement of the diffusion model's loss function to further improve synthetic data quality.
\end{abstract}

\begin{IEEEkeywords}
ion mobility spectrometry, synthetic data, diffusion models, data augmentation, classification
\end{IEEEkeywords}

\section{Introduction}
Ion mobility spectrometry (IMS) is widely used for rapid chemical detection in security, environmental monitoring, and defense applications, but obtaining large, balanced labeled datasets remains expensive and operationally constrained \cite{hill1990ion, dodds2019ion}. With only 1--50 labeled samples per class, traditional classifier training fails to capture intra-class variation and decision boundaries. We investigate whether synthetic data can address this limitation while preserving discriminative structure.

Prior work \cite{metzler2026, dunham2024oracle, preuer2019frechet, ho2020denoising} introduced a decoupled autoencoder with ChemNet-informed latent representations and our geometry-preserving diffusion process for synthetic spectrum generation. Here, we evaluate whether this synthetic data improves low-data classification.

Our preliminary investigation finds that synthetic augmentation provides substantial benefits for MLPs (+0.15aa to +0.21aa) and consistent benefits for RFs (+0.006aa to +0.01aa), with performance depending strongly on synthetic sampling variance. These results represent initial findings using default classifier hyperparameters. Comprehensive tuning and additional classifier architectures are under investigation.

\subsection{Contributions}
Our preliminary experiments show that geometry‑preserving synthetic data can meaningfully improve classifier performance when training data are scarce. Across a wide range of data budgets and mixing ratios, we observe consistent gains over real‑only training, and the operating‑point analysis highlights where synthetic samples provide the most benefit. These results were obtained using intentionally simple settings, and they establish a baseline that we will build on through more systematic tuning and model refinement in ongoing work.
%We present preliminary results on a controlled low-data benchmark spanning 50 data budgets, 11 mixing ratios, 2 classifiers, and 3 synthetic sampling scales (3,300 training configurations total). The evaluation provides quantitative evidence that synthetic augmentation produces large neural-classifier gains in low-data settings, and includes operating-point analysis identifying where synthetic variance and mixing ratio provide the greatest benefit. \textbf{FIXME REWRITE THIS TO NOT BE THE SAME} Ongoing work includes hyperparameter optimization for both classifiers and the generative model, evaluation of additional classifier architectures, and refinement of the diffusion model’s loss function to further improve synthetic data quality.
%Our preliminary results demonstrate that geometry-preserving synthetic data substantially improves neural classifier accuracy in low-data regimes, with operating-point analysis revealing optimal real/synthetic mixing ratios as a function of data budget. These baseline results establish clear performance gains using deliberately untuned configurations, providing a foundation for systematic optimization in ongoing work.

\section{Experimental Design}
\subsection{Data and Generative Source}
We use 8 chemical compounds from our IMS dataset \cite{dunham2024oracle}: DEB, DEM, DMMP, DPM, DtBP, JP8, MES, and TEPO. The base dataset contains 222,519 training and 74,173 test spectra. Synthetic spectra are generated from our latent diffusion pipeline operating in a 512-dimensional latent space at three sampling scales: $\sigma \in \{1.0, 1.5, 2.0\}$, where $\sigma$ controls the noise variance during sampling. For each scale, we generate 3,500 synthetic samples per class (84,000 total), then evaluate classifier performance under constrained data budgets by systematically varying real/synthetic mixtures.

\subsection{Classifier Evaluation Protocol}
For each $\sigma$ and ratio of real and synthetic data, we evaluate two classifiers:
\begin{itemize}
    \item RF (100 trees, default feature handling)
    \item MLP (architecture: 1000-500-250-100-8, adaptive Adam, standardized inputs)
\end{itemize}

For each class, we vary training size from 1 to 50 points/class and vary the real-data fraction $r \in \{0.0, 0.1, \dots, 1.0\}$. Here, $r=1.0$ is real-only training and $r=0.0$ is synthetic-only training. For each regime (fixed points/class), we compare:
\begin{equation}
\Delta = \max_{r \in [0,1]} \text{Acc}(r) - \text{Acc}(r=1.0)
\end{equation}
where positive $\Delta$ indicates that adding synthetic data beats real-only training.

\subsection{Reproducibility and Evaluation Scale}
All results in this paper come from the project evaluation pipeline\footnote{Code available at \url{https://github.com/metzlerk/Semi-Implicit-Invertibility-for-Data-Generation/tree/data_generation_essentials}.}. Each configuration is evaluated on the full 74,173-sample test set to ensure stable accuracy estimates. We report absolute accuracy deltas (not relative percent gain), and all classifier settings are intentionally untuned to isolate augmentation effects. The full experiment spans 3,300 unique training configurations (50 budgets $\times$ 11 mixing ratios $\times$ 2 classifiers $\times$ 3 sampling scales).

\section{Results}
\subsection{Overall Benefit Versus Real-Only Baseline}
Table~\ref{tab:summary} summarizes 50 low-data regimes per classifier and synthetic setting.

\begin{table*}[t]
\caption{Synthetic augmentation benefit relative to real-only training (50 regimes per row). Mean and Median $\Delta$ are in absolute accuracy (aa).}
\label{tab:summary}
\centering
\begin{tabular}{llcccc}
\toprule
\textbf{Classifier} & \textbf{Synthetic Std} & \textbf{Mean $\Delta$ (aa)} & \textbf{Median $\Delta$ (aa)} & \textbf{Improved Regimes} & \textbf{Synthetic-Only Better} \\
\midrule
RF & 1.0 & +0.0097 & +0.0000 & 20/50 & 2/50 \\
RF & 1.5 & +0.0061 & +0.0017 & 30/50 & 0/50 \\
RF & 2.0 & +0.0080 & +0.0006 & 28/50 & 1/50 \\
\midrule
MLP & 1.0 & +0.0129 & +0.0011 & 25/50 & 2/50 \\
MLP & 1.5 & +0.1519 & +0.1618 & 49/50 & 46/50 \\
MLP & 2.0 & +0.1304 & +0.1214 & 50/50 & 31/50 \\
\bottomrule
\end{tabular}
\end{table*}

\begin{figure*}[ht]
\centering
\includegraphics[width=\textwidth]{asilomar_2026_classifier_paper/Figure1.png}
\caption{Per-budget improvement over real-only training, defined as $\max_r \mathrm{Acc}(r) - \mathrm{Acc}(r{=}1.0)$ for each points-per-class budget. MLP gains are consistently large for std=1.5 and std=2.0, while RF gains are positive but comparatively small. (1 and 2 Points per class excluded)}
\label{fig:gain-over-real}
\end{figure*}

RF gains are present but small, roughly +0.006 to +0.01aa on average. This suggests that synthetic data helps coverage, but RFs are comparatively less sensitive to manifold expansion. In contrast, MLP benefits are substantial at higher synthetic variance. At std=1.5, the average best-minus-real gain is +0.1519aa, and real-only is surpassed in 49/50 regimes. At std=2.0, gains remain large at +0.1304aa on average, with improvement in all 50/50 regimes.

Figure~\ref{fig:gain-over-real} adds budget-level granularity: MLP gains are strongest in the low-data regime and taper as more real data is available, while RFs show smaller but stable positive gains.

\subsection{Budget-Stratified Analysis}
To better characterize where augmentation helps, Table~\ref{tab:budget-bins} reports mean gains in three budget bins.

\begin{table*}[ht]
\caption{Mean best-minus-real gain ($\Delta$ in absolute accuracy) by data budget bin.}
\label{tab:budget-bins}
\centering
\begin{tabular}{llccc}
\toprule
\textbf{Classifier} & \textbf{Synthetic Std} & \textbf{$n \leq 10$ (aa)} & \textbf{$11 \leq n \leq 30$ (aa)} & \textbf{$31 \leq n \leq 50$ (aa)} \\
\midrule
RF & 1.0 & +0.0415 & +0.0030 & +0.0005 \\
RF & 1.5 & +0.0120 & +0.0038 & +0.0054 \\
RF & 2.0 & +0.0260 & +0.0049 & +0.0021 \\
\midrule
MLP & 1.0 & +0.0325 & +0.0059 & +0.0101 \\
MLP & 1.5 & +0.2144 & +0.1978 & +0.0748 \\
MLP & 2.0 & +0.2144 & +0.1630 & +0.0558 \\
\bottomrule
\end{tabular}
\end{table*}

The budget-stratified view shows that synthetic augmentation provides the largest marginal value when real data is scarcest, particularly for MLPs. Even in higher-data bins, MLPs retain non-trivial gains for std=1.5 and std=2.0.

\subsection{Interpretation}
MLPs' higher capacity allows them to exploit manifold-supported variation from synthetic samples. Tree ensembles capture low-data signal with smaller marginal benefit from synthetic diversity.

At std=1.5, pure synthetic training (no real data) outperforms pure real training in 46/50 MLP regimes. The diffusion model has learned generalizable discriminative structure beyond memorization. The decline in gains at std=2.0 shows that excessive sampling variance introduces distribution shift that partially offsets augmentation benefits.

Across both models, std=1.5 provides the best trade-off: sufficient spread to improve class support and prevent overfitting, without excessive distribution shift.

\begin{figure*}[ht]
\centering
\includegraphics[width=\textwidth]{asilomar_2026_classifier_paper/Figure2.png}
\caption{Optimal real-data fraction $r^*$ (the ratio that maximizes accuracy for each budget). For MLPs at std=1.5 and std=2.0, optimal training often uses substantial synthetic content in low-data settings, while RFs remain closer to mixed or real-leaning ratios. (1 and 2 Points per class excluded)}
\label{fig:optimal-ratio}
\end{figure*}

\section{Discussion}
Each classifier benefits differently from synthetic augmentation. For MLPs in low-data regimes, gains are substantial: at std=1.5 with $n \leq 10$ real points/class, synthetic data improves accuracy by +0.2144aa on average. RFs show smaller but consistent gains (mean: +0.006aa to +0.010aa, peaking at $n \leq 10$ with +0.0415aa at std=1.0). Sampling scale matters: std=1.5 balances augmentation benefit against distribution shift risk.

Figure~\ref{fig:optimal-ratio} shows how the optimal real/synthetic ratio shifts with data budget. For neural models, synthetic-heavy or synthetic-only training is often optimal at low $n$, while real-heavy mixing becomes preferable as $n$ increases.

Evaluating generative models requires downstream classifier lift, not only moment matching or visual plausibility. Our geometry-preserving latent structure is critical: without manifold awareness, synthetic data adds noise rather than signal.

\section{Ongoing Development Leading to Full Paper}
These results represent preliminary findings using untuned classifier configurations. Several directions are under development that would be included in a full paper.

Classifier hyperparameters have not been systematically optimized. The RFs (100 trees, default settings) and MLPs (architecture: 1000-500-250-100-8, adaptive Adam) may benefit from grid search or Bayesian optimization over key parameters (tree depth, neurons, learning rate, dropout), which could reveal additional gains or alter the relative benefit of synthetic augmentation.

As detailed in our companion ICMLA submission \cite{metzler2026}, our latent diffusion model uses combined noise prediction and centroid-separation loss (80\%/20\% weighting). Ongoing work investigates alternative weighting schemes, additional geometric regularizers (Wasserstein distance preservation, local manifold alignment), and adaptive margin selection to improve synthetic sample quality.

Current evaluation uses fixed sampling scales $\sigma \in \{1.0, 1.5, 2.0\}$. We are investigating adaptive variance schedules that adjust diversity based on real data budget, class-specific variances for different chemical manifold structures, and curriculum approaches that progressively increase generation diversity during training.

\section{Conclusion}
This extended summary demonstrates that our geometry-preserving latent diffusion can materially improve low-data IMS spectra classification. Using baseline classifiers, synthetic augmentation yields large gains for neural networks (+0.15aa to +0.21aa at appropriate variance) and modest improvements for RFs (+0.006aa to +0.010aa).

These results establish proof-of-concept for synthetic augmentation in low-data IMS settings. As described in Section~V, comprehensive hyperparameter tuning, diffusion model refinements, and adaptive variance scheduling are under development. The full paper will incorporate these ongoing investigations.

All reported results use deliberately untuned configurations to isolate synthetic augmentation effects. The magnitude of gains with default settings suggests substantial room for improvement through systematic optimization.

\section*{Acknowledgments}
The authors would like to thank Casandra W. Maier and Zahra Bolkheiri for their valuable contributions to this work.

\bibliographystyle{IEEEtran}
\bibliography{bibliography}

\end{document}
